\documentclass[11pt,a4paper]{article}

\usepackage[T1]{fontenc}
\usepackage[utf8]{inputenc}
\usepackage[french]{babel}
\usepackage{lmodern}
\usepackage{microtype}
\usepackage[a4paper,margin=2.4cm]{geometry}
\usepackage{enumitem}
\usepackage[autostyle]{csquotes}
\usepackage[hidelinks]{hyperref}
\usepackage[backend=biber,style=authoryear,maxcitenames=2,maxbibnames=12]{biblatex}
\addbibresource{versions/references-v2.bib}

\setlength{\parindent}{0pt}
\setlength{\parskip}{0.65em}
\setlist[itemize]{leftmargin=1.5em,itemsep=0.35em,topsep=0.35em}
\setlist[enumerate]{leftmargin=1.8em,itemsep=0.5em,topsep=0.4em}

\newcommand{\projectversion}{2}
\newcommand{\projectdate}{12 juillet 2026}

\title{L'émergence de l'intéressant\\
\large Naissance des idées, compréhension et singularité des formes}
\author{François Pachet}
\date{Projet de thèse -- version \projectversion\\\small \projectdate}

\begin{document}

\maketitle

\begin{abstract}
Ce projet vise à constituer l'intéressant comme objet philosophique. Il part
d'une expérience ordinaire mais difficile à stabiliser : une idée surgit, une
forme devient intelligible, un passage retient l'attention et appelle la reprise.
La thèse centrale est que ces événements ne relèvent ni d'une propriété isolée
de l'objet, ni d'une préférence privée, mais d'une relation temporelle entre une
forme organisée et un sujet doté d'une histoire perceptive. La musique fournit
le laboratoire principal de cette enquête ; quarante années de construction de
systèmes d'intelligence artificielle en constituent un second terrain, à la fois
expérimental et réflexif. L'enjeu est d'articuler phénoménologie, philosophie des
formes, esthétique, sciences cognitives et modélisation computationnelle sans
réduire l'expérience de l'intéressant à l'un de ces registres.
\end{abstract}

\section{Origine et objectif du projet}

Ce travail est issu d'une pratique constructiviste et empirique de la recherche
en intelligence artificielle. J'ai passé quarante ans à construire des systèmes
pour générer, transformer ou accompagner des musiques et des textes. Le projet
doctoral ne consiste pas à résumer ces travaux ni à leur donner après coup une
unité philosophique artificielle. Il part d'une question plus exigeante :
\emph{comprendre ce que je fais} lorsque je construis un système, reconnais une
solution, élimine une production correcte mais sans intérêt, ou m'arrête devant
une forme qui semble s'imposer.

Cette enquête rencontre des problèmes que la philosophie a traités sous des noms
différents : accès à l'intelligible chez Platon et Aristote
\parencite{platon2016republique,platon1995theetete,aristote1991metaphysique,aristote2005seconds};
intuition et évidence chez Descartes et Spinoza
\parencite{descartes2002regles,descartes1992meditations,spinoza1966ethique};
jugement chez Kant \parencite{kant2001raison,kant1993juger}; invention et
reconfiguration des problèmes chez Peirce, Poincaré, Hadamard, Bachelard et
Valéry \parencite{peirce1955writings,poincare1908science,hadamard1959invention,bachelard1999formation,valery1973cahiers};
mémoire, durée et constitution temporelle du sens chez Bergson, Husserl et
Merleau-Ponty
\parencite{bergson2013evolution,bergson2013pensee,husserl1970experience,husserl1957logique,merleau1945phenomenologie};
changement d'aspect chez Wittgenstein \parencite{wittgenstein2004recherches};
enfin symbolisation, différence et répétition chez Goodman et Deleuze
\parencite{goodman1976languages,deleuze1968difference}.

L'hypothèse d'unification est que la naissance d'une idée, l'expérience de la
compréhension et la fascination pour une forme sont trois moments d'un même
phénomène : l'émergence de quelque chose qui se détache pour un sujet comme une
organisation à la fois intelligible et singulière.

\section{Question directrice et thèse centrale}

La question directrice est : \emph{sous quelles conditions quelque chose
apparaît-il à un sujet comme intéressant ?} Une chanson, une démonstration, une
phrase ou une idée peut retenir l'attention de manière disproportionnée, provoquer
des reprises et des variations, puis perdre ou transformer son pouvoir. Cet effet
ne se réduit ni à la nouveauté brute, car l'aléatoire surprend sans nécessairement
intéresser, ni à la beauté, car une forme familière peut demeurer belle tout en
étant devenue prévisible, ni à la popularité, ni à une préférence subjective
entendue comme donnée sans histoire.

La thèse centrale est la suivante :

\begin{quote}
L'intéressant est une propriété relationnelle et temporelle. Il émerge lorsqu'une
forme organisée transforme l'état d'attention, de mémoire ou de compréhension
d'un sujet historiquement, affectivement, corporellement et culturellement
constitué.
\end{quote}

La relativité au sujet n'implique pas que tous les jugements se valent. Une
relation peut être décrite : on peut identifier les attentes rendues possibles
par l'objet, les compétences nécessaires pour les éprouver, les transformations
qui relancent l'attention, ainsi que les variantes qui détruiraient l'effet. La
structure de l'objet et l'histoire du sujet sont donc deux conditions
irréductibles de la même rencontre.

\section{Hypothèses de travail}

\begin{enumerate}
  \item \textbf{L'intéressant est un processus, non une étiquette.} Pour un objet
  temporel, l'unité pertinente est une trajectoire faite de préparation,
  d'attente, d'écart, de résolution, de réinterprétation et parfois d'ennui. Un
  même événement change de valeur selon sa place et selon ce que la suite fait de
  lui.

  \item \textbf{Une forme intéressante peut inventer le problème qu'elle
  résout.} Certaines solutions paraissent nécessaires après leur apparition,
  alors qu'elles n'étaient pas déductibles auparavant. Elles rendent visibles
  leurs propres contraintes et donnent l'impression d'une simplicité de surface
  qui dissimule une genèse improbable. L'intérêt pertinent tient alors moins à
  une rareté statistique brute qu'à la petitesse d'une classe de solutions dans
  un problème implicite.

  \item \textbf{L'intérêt se déplace avec l'apprentissage.} Une régularité
  d'abord opaque peut devenir accessible, produire un gain de compréhension,
  puis s'épuiser lorsqu'elle est acquise. La théorie du progrès de compression
  formalise une version de cette dynamique
  \parencite{schmidhuber1997interesting,schmidhuber2009compression}. Elle doit
  toutefois être comprise relativement aux capacités de calcul, à la mémoire et
  à la culture d'un observateur, et non comme une mesure universelle de l'objet.

  \item \textbf{L'intéressant se situe souvent à une frontière mouvante du
  presque-apprenable.} Le déjà connu est trivial ; le bruit demeure
  incompressible ; une structure trop éloignée des capacités du sujet reste
  opaque. Une forme intéresse lorsqu'elle promet une transformation accessible
  au prix d'un effort non nul. Cette proposition prolonge, tout en la compliquant,
  la courbe en U inversé de l'esthétique expérimentale
  \parencite{berlyne1974studies,abdallah2009information}.

  \item \textbf{L'attention est active et faillible.} L'auditeur varie,
  reconstruit et compare intérieurement ce qu'il entend. Cette activité peut
  révéler une organisation réelle, mais aussi fabriquer l'intérêt qu'elle croit
  découvrir. Familiarité, prestige de l'auteur et désir de résoudre une énigme
  sont des sources possibles de faux positifs. Une théorie de l'intéressant doit
  donc contenir sa propre critique.
\end{enumerate}

\section{Cadre théorique : sujet, forme et émergence}

Le cadre théorique cherchera à articuler trois exigences. La première vient de la
philosophie du sujet : l'expérience de l'intéressant dépend de conditions de
réception et de jugement, sans être réductible à un trait objectif. La deuxième
vient d'une phénoménologie de la temporalité : rétention, protention, mémoire
corporelle et changement d'aspect décrivent comment une forme se constitue dans
le temps. La troisième vient des théories morphodynamiques et des systèmes
complexes : une organisation globale possède des propriétés qui ne se lisent pas
dans ses éléments isolés.

Ce cadre reprend la tentative de naturaliser certains problèmes
transcendantaux sans éliminer leur dimension phénoménale
\parencite{petitot1993phenomenologie,petitot2002naturaliser}. Il ne postule pas
que la modélisation scientifique épuisera l'intéressant. Il cherche plutôt un
rapport contrôlé entre descriptions à la première personne, analyses formelles
et mécanismes testables. Il faudra distinguer les invariants possibles de
l'attention et les habitudes propres à un contexte, notamment celles de la
musique tonale occidentale.

\section{La musique comme laboratoire philosophique}

La musique rend particulièrement visibles les rapports entre mémoire, attente,
surprise et résolution. Meyer puis Narmour ont montré que l'affect et la
signification musicale dépendent d'attentes apprises, réalisées, retardées ou
déjouées \parencite{meyer1956emotion,narmour1990basic,narmour1992complexity}.
Les modèles informationnels et neuro-symboliques donnent des moyens de préciser
ces dynamiques
\parencite{abdallah2009information,berger1999expectations,gang1999unified}.

Mais expliquer pourquoi un titre est intéressant ne consiste pas à repérer une
quantité moyenne de surprise. L'enquête portera sur la trajectoire d'attention à
plusieurs échelles : micro-émotions locales, relations entre mélodie et harmonie,
forme globale, histoire des écoutes et activité de l'auditeur. Le régime Jotney
constitue un premier cas : autonomie mélodique et mobilité harmonique y créent
une bande étroite dans laquelle la surprise devient rétrospectivement nécessaire.

\emph{Histoire d'une oreille} fournit un vocabulaire phénoménologique et une
méthode \parencite{pachet2018oreille}. L'\emph{escatte}, l'\emph{escassepte}, la
\emph{misette}, la \emph{neuvéclamine}, la \emph{cébimine} ou le
\emph{stabilème} ne sont pas de nouveaux noms pour des accords ; ils décrivent
des transformations attentionnelles produites par des relations temporelles.
L'\emph{écoute superposée} soumet la musique à des variations intérieures.
L'\emph{optale} désigne le sentiment qu'une forme est un optimum local parce que
ses variantes proches échouent ; la \emph{pseudoptale} avertit que la familiarité
peut simuler cette nécessité. Les \emph{fausses évidences} distinguent une
banalité requalifiée après coup d'un simple \emph{substrat nécessaire}. Enfin,
l'\emph{écoute surattentive} nomme le risque que l'auditeur produise lui-même la
résolution qu'il attribue à l'œuvre.

Ces concepts permettent de passer de « ce titre me plaît » à des propositions
comparables : quelle transformation se produit, à quel moment, grâce à quelle
préparation, pour quelle oreille, et que deviendrait-elle sous une variation
contrôlée ? Ils n'ont pas vocation à former un dictionnaire définitif. Ils sont
des instruments pour décrire et mettre à l'épreuve une théorie de l'attention
musicale.

\section{Pratique scientifique et enquête philosophique}

Les systèmes d'IA constituent un second laboratoire. Les travaux sur Flow
Machines, les interactions réflexives, la génération markovienne sous contraintes
et l'\emph{interestingness} computationnelle ont produit des objets, des
protocoles, des réussites et des échecs
\parencite{pachet2006clefs,pachet2008machines,pachet2012virtuosity,pachet2021assisted,pachet2026markov}.
Ils montrent notamment qu'une génération locale en marche avant ne suffit pas
toujours à satisfaire une organisation globale : produire une forme intéressante
exige souvent d'articuler échantillonnage et résolution de contraintes.

Trois niveaux devront rester distincts :

\begin{itemize}
  \item les résultats scientifiques et les propriétés effectives des systèmes ;
  \item les propositions philosophiques sur l'invention, l'attention ou la
        valeur ;
  \item les articulations argumentées qui expliquent ce que la construction d'un
        système rend pensable sans prétendre qu'elle le démontre.
\end{itemize}

Cette articulation est elle-même une contribution attendue de la thèse. Un
modèle computationnel peut rendre une hypothèse explicite, produire un
contre-exemple ou révéler qu'une notion est trop générale. Il ne suffit pas à
naturaliser l'expérience qu'il schématise.

\section{Méthode et corpus}

La recherche combinera :

\begin{itemize}
  \item une analyse conceptuelle de l'intéressant et de ses voisins : nouveau,
        surprenant, beau, pertinent, difficile, rare et compréhensible ;
  \item une lecture des traditions philosophiques de la forme, du jugement, de
        la temporalité, de l'invention et de l'émergence ;
  \item une phénoménologie située de cas musicaux, littéraires, mathématiques et
        scientifiques, incluant l'histoire de la formation d'une oreille ;
  \item des comparaisons contrefactuelles : variations intérieures,
        continuations simplement compétentes, reprises ou imitations ratées ;
  \item l'analyse réflexive de systèmes construits et, lorsque c'est pertinent,
        de nouvelles expériences sur l'attention, la reprise, l'abandon et
        l'évolution des jugements ;
  \item une confrontation avec les modèles contemporains de l'attente, de la
        compression, de la curiosité et de la génération contrainte.
\end{itemize}

Le corpus principal réunit les deux essais sur l'impossibilité de créer et la
virtuosité \parencite{pachet2026impossibilite,pachet2026virtuosite},
\emph{Histoire d'une oreille}, les travaux sur les chansons Jotney, les projets
scientifiques sur l'\emph{interestingness} et les systèmes génératifs, ainsi que
les sources philosophiques et scientifiques citées dans la bibliographie. Les
cas anciens ne seront retenus que lorsqu'ils éclairent directement la question
de l'intéressant.

\section{Contributions attendues}

La thèse vise cinq contributions :

\begin{enumerate}
  \item une définition relationnelle et temporelle de l'intéressant, distincte
        de la nouveauté, de la beauté et de la préférence ;
  \item une théorie de la nécessité inventée, reliant forme, problème implicite,
        rareté des solutions et genèse improbable ;
  \item une théorie de l'attention musicale capable d'articuler micro-émotions,
        attentes, forme globale, apprentissage et activité de l'auditeur ;
  \item une méthode de passage entre construction scientifique et proposition
        philosophique, avec des critères explicites de portée et de limite ;
  \item un ensemble d'hypothèses falsifiables ou au moins comparables sur les
        trajectoires d'intérêt, leurs transformations et leurs faux positifs.
\end{enumerate}

\section{Plan provisoire}

\begin{enumerate}
  \item \textbf{Distinguer l'intéressant} : concept, voisins, histoire du
        problème et objections.
  \item \textbf{Émergence et compréhension} : apparition d'une forme,
        changement d'aspect et naissance des problèmes.
  \item \textbf{Forme, mémoire et attention} : temporalité, apprentissage,
        presque-apprenable et épuisement de l'effet.
  \item \textbf{Nécessité inventée} : contraintes, rareté, autostabilité,
        difficulté de production et simplicité de surface.
  \item \textbf{Pourquoi ce titre est-il intéressant ?} : attentes musicales,
        micro-émotions, optales, forme globale et surattention.
  \item \textbf{Construire pour comprendre} : systèmes d'IA, expériences,
        portée philosophique et limites de la naturalisation.
\end{enumerate}

\section{Travaux personnels directement liés}

\begin{itemize}
  \item \textcite{pachet2026virtuosite} analyse les formes de virtuosité et leurs
        effets recherchés sur les auditeurs.
  \item \textcite{pachet2026impossibilite} examine les obstacles scientifiques
        et conceptuels à l'idée de créer du nouveau.
  \item \textcite{pachet2026markov} synthétise les travaux sur la modélisation
        markovienne et la génération contrôlée par contraintes.
  \item \textcite{pachet2021assisted} tire des projets Flow Machines de nouvelles
        catégories du nouveau.
  \item \textcite{pachet2018oreille} décrit l'ontogenèse d'une oreille musicale
        et fournit les concepts attentionnels mobilisés ici.
  \item \textcite{pachet2012virtuosity} relie virtuosité, créativité et
        construction du système Virtuoso.
  \item \textcite{pachet2008machines} présente les systèmes réflexifs comme des
        machines favorisant la sortie de soi.
  \item \textcite{pachet2006clefs} propose une première modélisation des mélodies
        intéressantes.
\end{itemize}

\printbibliography[title={Références}]

\end{document}
